\slajdid{09_1-s003}
\begin{frame}[label=09_1-s003]{NMOS indukovan „dugački“ kanal}
\gusttekst\setlength{\parskip}{2pt}\setlength{\abovedisplayskip}{3pt}\setlength{\belowdisplayskip}{3pt}\setlength{\abovedisplayshortskip}{2pt}\setlength{\belowdisplayshortskip}{2pt}
\begin{columns}[T,onlytextwidth]\column{.49\textwidth}\centering
\begin{tikzpicture}[x=.8cm,y=.65cm,font=\sitneoznake,>=Latex]
\foreach\x in{0,3}{\begin{scope}[shift={(\x,0)}]
\draw(-.6,0)--(-.1,0);\draw(-.1,-.35)--(-.1,.35);\draw(0,-.4)--(0,.4);\draw(0,.3)--(.4,.3)--(.4,.9)node[above]{D};\draw(0,-.3)--(.4,-.3)--(.4,-.9)node[below]{S};\ifnum\x=3\node[left]at(-.9,0){G};\else\node[left]at(-.6,0){G};\fi\end{scope}}
\draw[->](-.7,.2)--(-.2,.2)node[midway,above]{$I_G$};\draw[->](.65,.85)--(.65,.35)node[midway,right]{$I_D$};\draw[->](.65,-.35)--(.65,-.85)node[midway,right]{$I_S$};
\draw[<->](3.4,.85)to[bend right=40](2.35,.05);\node at(2.5,.95){$V_{GD}$};\node at(2.25,.24){$+$};
\draw[<->](2.35,-.1)to[bend right=40](3.4,-.85);\node at(2.5,-.9){$V_{GS}$};\node at(2.2,-.22){$+$};
\draw[<->](3.65,.8)to[bend left=20](3.65,-.8);\node[right]at(3.85,0){$V_{DS}$};\node at(3.7,.95){$+$};
\end{tikzpicture}

\[I_{Gn}=0\qquad I_{Dn}=I_{Sn}\geq0\quad za\ V_{GSn}\geq V_{Tn}\]
\raggedright Tranzistor vodi u zasićenju, aktivnoj oblasti, kada je
\[V_{DSn}\geq V_{GSn}-V_{Tn}\]
i tada je njegova struja data izrazom
\[I_{Dn}=\frac{k_n}{2}(V_{GSn}-V_{Tn})^2(1+\lambda_nV_{DSn})\]
Tranzistor vodi u triodnoj, omskoj, neaktivnoj oblasti, kada je
\[V_{DSn}<V_{GSn}-V_{Tn}\]
i tada je njegova struja data izrazom
\[I_{Dn}=\frac{k_n}{2}\bigl(2V_{DSn}(V_{GSn}-V_{Tn})-V_{DSn}^2\bigr)\]
\column{.49\textwidth}
$V_{Tn}>0$ - prag uključenja n kanalnog tranzistora, napon pri kojem se formira kanal\par\medskip
$\lambda_n$ – parametar koji pokazuje modifikaciju dužine kanala u zavisnosti od napona između drejna i sors – jedinica $\frac1V$
\[k_n=\mu_nC_{oxn}\frac{W_n}{L_n}\]
$W_n$ – širina kanala\par\medskip
$L_n$ – dužina kanala\par\medskip
$\mu_n$ – pokretljivost nosilaca koja redstavlja vezu između brzine nosilaca i polja koje deluje na njih $v=\mu E$ - jedinica $\frac{m^2}{Vs}$\par\medskip
$C_{oxn}$ – normalizovana, po površini, kapacitivnost gejta– jedinica $\frac F{m^2}$\par\medskip
$k_n$ – jedinica $\frac A{V^2}$
\end{columns}
\end{frame}
